% Գլուխ 1. Տեսական հիմքեր (սևագիր v1)

\chapter{Տեսական հիմքեր}

\section{Գրաֆների հիմնական հասկացություններ}

\begin{definition}
\emph{Պարզ չուղղորդված գրաֆ} $G=(V,E)$ անվանում են զույգ, որտեղ $V$-ն վերջավոր բազմություն է (\emph{գագաթներ}), իսկ $E$-ն $V$-ի երկտարր ենթաբազմությունների բազմություն է (\emph{կողեր}): Նշանակենք $n=|V|$, $m=|E|$:
\end{definition}

Գագաթի $v$ \emph{աստիճան} $\deg(v)$ կոչվում է նրան կից կողերի թիվը: Գրաֆի նվազագույն աստիճանը նշանակենք $\delta(G)$:

\begin{definition}
\emph{Ուղի} է կոչվում գագաթների $v_0,v_1,\dots,v_k$ հաջորդականությունը, որտեղ ամեն $\{v_{i-1},v_i\}$ կող է և գագաթները զույգ առ զույգ տարբեր են: \emph{Ցիկլ} է կոչվում փակ ուղին (երբ $k\ge 3$ և $v_k=v_0$, մնացած գագաթները տարբեր են):
\end{definition}

\begin{definition}
Գրաֆը \emph{կապակցված} է, եթե նրա ցանկացած երկու գագաթների միջև գոյություն ունի ուղի: Գրաֆի \emph{կապակցվածության բաղադրիչը} նրա մաքսիմալ կապակցված ենթագրաֆն է:
\end{definition}

Գրաֆը կապակցված է այն և միայն այն դեպքում, երբ ունի ուղիղ մեկ կապակցվածության բաղադրիչ: Բաղադրիչները կարելի է գտնել BFS կամ DFS ալգորիթմով $O(n+m)$ ժամանակում (տե՛ս Գլուխ 2):

\begin{example}
Ցիկլը $C_n$ կապակցված է: Եթե նրանից հեռացնենք ցանկացած մեկ կող, կստացվի ուղի, որը դեռ կապակցված է: Եթե հեռացնենք ցանկացած երկու կող, գրաֆը կտրոհվի ուղիղ երկու կապակցվածության բաղադրիչի:
\end{example}

\section{Կապակցվածություն, կտրող կետեր, կամուրջներ}

\begin{definition}
Կապակցված գրաֆի $v$ գագաթը կոչվում է \emph{կտրող կետ}, եթե $G-v$ գրաֆը (գագաթը և նրան կից կողերը հեռացնելուց հետո) կապակցված չէ: Կապակցված գրաֆի $e$ կողը կոչվում է \emph{կամուրջ}, եթե $G-e$ գրաֆը կապակցված չէ:
\end{definition}

Ցանցային տերմիններով կտրող կետը \emph{միակ խոցելի հանգույցն} է, իսկ կամուրջը՝ \emph{միակ խոցելի կապը}, որի խափանումը բաժանում է ցանցը մասերի:

\begin{theorem}\label{thm:bridge}
Կապակցված $G$ գրաֆի $e=\{u,v\}$ կողը կամուրջ է այն և միայն այն դեպքում, երբ $e$-ն չի պատկանում ոչ մի ցիկլի:
\end{theorem}

\begin{proof}
Ենթադրենք $e$-ն պատկանում է ինչ-որ ցիկլի $C$: Այդ դեպքում $G-e$ գրաֆում $u$-ից $v$ ուղի կա՝ ցիկլի մնացած մասը: Ցանկացած ուղի, որն անցնում էր $e$-ով, կարելի է շրջանցել այդ ուղիով, ուստի $G-e$-ն կապակցված է, և $e$-ն կամուրջ չէ:

Հակառակը, ենթադրենք $e$-ն ոչ մի ցիկլում չի գտնվում, բայց $G-e$ կապակցված է: Այդ դեպքում $G-e$-ում կա $u$-ից $v$ ուղի $P$, և $P$-ն $e$-ի հետ կազմում է ցիկլ, որը պարունակում է $e$: Սա հակասում է ենթադրությանը, ուստի $G-e$ կապակցված չէ, այսինքն՝ $e$-ն կամուրջ է:
\end{proof}


\section{\texorpdfstring{$k$}{k}-կապակցվածություն և Մենգերի թեորեմը}

\begin{definition}
$G$ գրաֆի \emph{գագաթային կապակցվածություն} $\kappa(G)$ կոչվում է գագաթների նվազագույն թիվը, որոնց հեռացումը գրաֆը դարձնում է ոչ կապակցված կամ մեկ գագաթանի (լրիվ գրաֆի համար $\kappa(K_n)=n-1$): \emph{Կողային կապակցվածություն} $\lambda(G)$ կոչվում է կողերի նվազագույն թիվը, որոնց հեռացումը գրաֆը դարձնում է ոչ կապակցված: Գրաֆը \emph{$k$-կապակցված} է (կողային $k$-կապակցված), եթե $\kappa(G)\ge k$ (համապատասխանաբար $\lambda(G)\ge k$):
\end{definition}

\begin{theorem}[Ուիթնի անհավասարություն]
Ցանկացած գրաֆի համար $\kappa(G)\le\lambda(G)\le\delta(G)$:
\end{theorem}

Թեորեմի ապացույցը տե՛ս \cite{west}: Դրա գործնական իմաստը. ցանցի կապի հուսալիությունը վերևից սահմանափակված է ամենափոքր աստիճան ունեցող հանգույցով:

Մենգերի թեորեմը կապում է կապակցվածությունը \emph{չհատվող ուղիների} թվի հետ: Սա կապակցվածության չափման ալգորիթմների հիմքն է (Գլուխ 2):

\begin{theorem}[Մենգեր, 1927]\label{thm:menger}
Դիցուք $s,t$ երկու ոչ հարևան գագաթներ են: $s$-ը և $t$-ն իրարից բաժանող գագաթների նվազագույն թիվը հավասար է $s$-ից $t$ գագաթային չհատվող (ներքին գագաթներով անհատ) ուղիների առավելագույն թվին: Կողային տարբերակում՝ $s$ և $t$ գագաթները բաժանող կողերի նվազագույն թիվը հավասար է $s$-ից $t$ կողային չհատվող ուղիների առավելագույն թվին:
\end{theorem}

Հետևանք. գրաֆը $k$-կապակցված է այն և միայն այն դեպքում, երբ ցանկացած երկու գագաթի միջև կան առնվազն $k$ գագաթային չհատվող ուղիներ (Ուիթնի թեորեմ, տե՛ս \cite{west}):

\begin{example}
Ցիկլ $C_n$ ($n\ge 3$)-ում ցանկացած երկու գագաթի միջև կան ուղիղ 2 գագաթային չհատվող ուղիներ (երկու ուղղություններով), ուստի $\kappa(C_n)=\lambda(C_n)=2$: Ծառի համար $\kappa=\lambda=1$, քանի որ ցանկացած կող կամուրջ է:
\end{example}

\section{Ցանցի հուսալիության մոդելներ}

Մինչ այժմ դիտարկվում էր ցանցի \emph{ստատիկ} կառուցվածքը: Հուսալիության վերլուծության մեջ ենթադրվում է, որ ցանցի տարրերը կարող են պատահականորեն խափանվել:

\begin{definition}
\emph{Հավանականային գրաֆ} է կոչվում $G=(V,E)$ գրաֆը, որի ամեն $e\in E$ կող աշխատում է $p_e$ հավանականությամբ և խափանվում $1-p_e$ հավանականությամբ, անկախ մյուս կողերից: Հաճախ ընդունում են $p_e=p$ բոլոր կողերի համար: Գագաթները համարվում են բացարձակ հուսալի:
\end{definition}

Ցանցի վիճակը որոշվում է աշխատող կողերի $S\subseteq E$ բազմությամբ, որի հավանականությունն է
\[
\Pr(S)=\prod_{e\in S}p_e\prod_{e\notin S}(1-p_e).
\]

\begin{definition}
\emph{Բոլոր գագաթների հուսալիություն} (all-terminal reliability) $R_{\mathrm{all}}(G)$ կոչվում է այն հավանականությունը, որ աշխատող կողերով ենթագրաֆը (բոլոր $V$ գագաթներով) կապակցված է.
\[
R_{\mathrm{all}}(G)=\sum_{S\subseteq E:\ (V,S)\text{ կապակցված}}\Pr(S).
\]
\emph{Երկու գագաթների հուսալիություն} (two-terminal reliability) $R_{s,t}(G)$ կոչվում է այն հավանականությունը, որ աշխատող կողերով ենթագրաֆում $s$-ից $t$ ուղի գոյություն ունի:
\end{definition}

\begin{example}\label{ex:reliab}
Դիցուք բոլոր կողերն ունեն նույն $p$ հավանականությունը:
\begin{itemize}
  \item Ծառ $n$ գագաթով ($n-1$ կող). բոլոր կողերը պետք է աշխատեն, ուստի $R_{\mathrm{all}}=p^{\,n-1}$:
  \item Ցիկլ $C_n$. ցանցը կապակցված է, երբ աշխատում են բոլոր $n$ կողերը կամ ուղիղ $n-1$-ը, ուստի
  \[
  R_{\mathrm{all}}(C_n)=p^{\,n}+n\,p^{\,n-1}(1-p).
  \]
  Օրինակ, $n=3$, $p=0{,}9$ դեպքում $R=0{,}729+3\cdot 0{,}81\cdot 0{,}1=0{,}972$:
  \item Երկու գագաթների միջև \emph{հաջորդական} կապ ($p_1,p_2$ հավանականությամբ երկու կող). $R_{s,t}=p_1p_2$: \emph{Զուգահեռ} կապ. $R_{s,t}=1-(1-p_1)(1-p_2)$:
\end{itemize}
\end{example}

Այս օրինակները ցույց են տալիս, որ պարզ կառուցվածքների համար հուսալիությունը հաշվվում է փակ բանաձևով: Ընդհանուր դեպքում, սակայն, խնդիրը դժվար է. ճշգրիտ հաշվարկի բարդությունը աճում է էքսպոնենցիալ ($2^m$ վիճակ), և հայտնի է, որ $R_{\mathrm{all}}$ և $R_{s,t}$ հաշվարկը \#P-դժվար է \cite{valiant}: Դա է պատճառը, որ գործնականում օգտագործում են մոտավոր մեթոդներ, մասնավորապես Monte Carlo սիմուլյացիա (տե՛ս Գլուխ 2):

% TODO: ավելացնել գագաթային խափանումներով մոդել, կապը կապակցվածության հետ (R >= p^lambda-ի տիպի գնահատականներ)
